\chapter{Introduction}\label{introduction}

\section{Motivation}
The digitization of healthcare has led to the widespread use of Electronic Health Records, which are inherently heterogeneous, high-dimensional, and longitudinal. These characteristics make clinical risk prediction particularly challenging~\cite{ossboll2024gnn}. In recent years, Graph Neural Networks have emerged as a promising approach in this domain, since representing patients as graphs of medical entities allows them to capture relational and temporal structures that traditional models often fail to model effectively~\cite{ossboll2024gnn,mesinovic2026dynagraph}.

Despite their strong predimctive performance, GNNs still suffer from an important limitation: they remain largely black box models. As noted in prior work, ``a major limitation of deep models is that they are not amenable to interpretability''~\cite{yuan2022explainability}. In clinical settings, where transparency is essential for trust and medical decision making, this lack of interpretability is especially problematic.

Most existing work addresses this issue using post-hoc explanation methods such as GNNExplainer~\cite{ying2019gnnexplainer}. However, these methods are applied after the model has already made its prediction, which makes it unclear whether the explanations truly reflect how the model arrived at its decision. In addition, their performance can be inconsistent when applied to larger datasets or more complex evaluation settings~\cite{agarwal2023graphxai}.

An alternative approach is using methods with built-in interpretability whose internal decision-making logic is transparent and directly explainable to humans. One of such methods is ProtGNN, which is a prototype-based model that explains its decisions by comparing a test instance to learned representative cases, driving predictions based on similarities ~\cite{zhang2022protgnn}. However, most work on this type of model has been tested on molecular data, so it is still unclear how well it works for healthcare applications. The closest related work in this area, ProtoEHR, applies a similar idea to Electronic Health Records data, but it does not identify which parts of the patient graph drive the prediction and does not compare its explanations with post-hoc methods in a consistent way~\cite{cai2025protoehr}.

This project aims to address this gap by developing a self-explainable GNN for EHR-based diagnosis that incorporates a built-in explainable layer. Three models were implemented on the same patient dataset, where each model is designed to make predictions and to extract the medical subgraph that drives each of those predictions, specifically through prototype matching, internal attention mechanisms, or 3RD METHOD'S APPROACH. Each approach will be evaluated in terms of both predictive performance and explanation quality, and compared against post-hoc baselines using the GraphXAI framework~\cite{agarwal2023graphxai}.

\section{Description of Work}
This Master's project is structured in three phases. First, a focused literature review will be conducted on GNN-based clinical prediction from EHR data, post-hoc explanation methods and their limitations, and different methods with built-in interpretability.

Second, an EHR-to-graph pipeline will be developed using the MIMIC-IV-ED dataset~\cite{johnson2023mimicived}. Medical entities such as diagnoses, medications, lab results, and procedures will be represented as nodes, while edges will capture temporal and clinical relationships. To achieve interpretability, built-in explainability layers will be embedded into the graph neural network architectures to identify clinical patterns in predictions and highlight the most relevant medical subgraphs.

Third, the model will be evaluated along two dimensions. Predictive performance will be measured with multi-class metrics and compared against a tabular baseline and other GNN methods trained on the same split. Explanation quality will be assessed with sparsity and inter-explainer agreement, and compared with post-hoc methods within the GraphXAI framework~\cite{agarwal2023graphxai}.

The project aims to provide a reproducible self-explainable GNN framework for EHR-based diagnosis, documented in a report aligned with scientific research standards.

\chapter{Project Outline}

This chapter should define the scope, objectives, research questions, methodology, expected contributions, and organization of the project.

\textbf{RQ1}: How does the proposed prototype-based self-explainable GNN perform compared with tabular, standard GNN, and knowledge-graph baselines for imbalanced multi-class diagnosis?

\textbf{RQ2}: What is the contribution of the prototype layer and prototype-projection mechanism to predictive performance and interpretability?

\textbf{RQ3}: How faithful, stable, sparse, and clinically meaningful are the intrinsic and post-hoc explanations generated by the proposed approach?

